% D26  -
% Auto-generated from markdown source
% origin: ai

\documentclass[12pt, a4paper]{article}
\usepackage[utf8]{inputenc}
\usepackage[T1]{fontenc}
\usepackage{amsmath, amssymb, amsfonts, amsthm}
\usepackage{textcomp}
\usepackage{geometry}
\usepackage{hyperref}
\usepackage{booktabs}
\usepackage{graphicx}
\usepackage{fancyhdr}
\usepackage{tcolorbox}
\usepackage{enumitem}

\geometry{margin=1in}
\pagestyle{fancy}
\fancyhf{}
\rhead{D26  -}
\lhead{\thepage}

\newtheorem{definition}{Definition}[section]
\newtheorem{theorem}{Theorem}[section]
\newtheorem{corollary}{Corollary}[section]
\newtheorem{proposition}{Proposition}[section]
\newtheorem{lemma}{Lemma}[section]
\newtheorem{conjecture}{Conjecture}[section]

\newcommand{\calI}{\mathcal{I}}
\newcommand{\calP}{\mathcal{P}}
\newcommand{\calE}{\mathcal{E}}
\newcommand{\calS}{\mathcal{S}}
\newcommand{\calH}{\mathcal{H}}
\newcommand{\calO}{\mathcal{O}}
\newcommand{\calF}{\mathcal{F}}
\newcommand{\calU}{\mathcal{U}}
\newcommand{\calV}{\mathcal{V}}
\newcommand{\calA}{\mathcal{A}}
\newcommand{\calG}{\mathcal{G}}
\newcommand{\calB}{\mathcal{B}}
\newcommand{\calC}{\mathcal{C}}
\newcommand{\calL}{\mathcal{L}}
\newcommand{\calM}{\mathcal{M}}
\newcommand{\calX}{\mathcal{X}}
\newcommand{\calY}{\mathcal{Y}}
\newcommand{\calW}{\mathcal{W}}
\newcommand{\calT}{\mathcal{T}}
\newcommand{\calZ}{\mathcal{Z}}
\newcommand{\calR}{\mathcal{R}}
\newcommand{\calQ}{\mathcal{Q}}
\newcommand{\calD}{\mathcal{D}}
\newcommand{\calN}{\mathcal{N}}
\newcommand{\calK}{\mathcal{K}}

\newcommand{\Tr}{\text{Tr}}
\newcommand{\E}{\mathbb{E}}
\newcommand{\diff}{\mathrm{d}}
\newcommand{\KL}{D_{\mathrm{KL}}}
\newcommand{\argmax}{\arg\max}
\newcommand{\argmin}{\arg\min}
\newcommand{\hbarF}{\hbar}

\newtcolorbox{thmbox}[1]{
  colback=blue!5!white, colframe=blue!75!black, fonttitle=\bfseries,
  title={#1}, breakable
}
\newtcolorbox{defbox}[1]{
  colback=green!5!white, colframe=green!75!black, fonttitle=\bfseries,
  title={#1}, breakable
}

\begin{document}

\title{D26  -}
\author{RALPH Loop}
\date{}
\maketitle

\begin{abstract}
This paper presents D26  -, formalizing the underlying principles using information-theoretic methods.

\textbf{Keywords}: D26  -, information theory, formalization.
\end{abstract}

\tableofcontents
\section{D26  -}




\subsection{0.}


****
\begin{itemize}
\item ****unstable
\item ****stable
\item ****critical
\end{itemize}


\textbf{D26 }\textbf{*}\textbf{*}**



\subsection{1.}


\subsubsection{1.1}


\textbf{ D26-1}
LTI
\[\dot{x} = Ax + Bu, \quad y = Cx + Du\]

****
\[\dot{x} = (A + \delta A)x + (B + \delta B)u + \eta\]


\begin{itemize}
\item \(\delta A, \delta B\)
\item \(\eta\)
\end{itemize}



\subsubsection{1.2}


\textbf{ D26-2 \(G_{AMP}\)}
\(\omega\)
\[G_{AMP}(\omega) = \frac{\|y(t; \delta)\|_2}{\|\delta\|_2}\]

****
\[G(s) = C(sI - A)^{-1}B + D\]

\textbf{ = }
\[G_{AMP}(\omega) = \|G(j\omega)\|_2 = \sigma_{\max}(G(j\omega))\]


\subsubsection{1.3}


\textbf{ D26-3}
****
\[\exists \lambda_i \in \text{spec}(A): \text{Re}(\lambda_i) = 0\]

****
\[\kappa_c = \min_{\lambda_i} |\text{Re}(\lambda_i)|\]


\subsection{2.}


\subsubsection{ D26-1}


****
\textbf{*}***
\[S_{\lambda_i}(A) = \frac{\partial \lambda_i}{\partial A} = v_i w_i^\top\]

\(v_i, w_i\) /

\textbf{D26-1 }
\[\lim_{t \to \infty} \|x(t)\| = \infty \iff \exists i: \text{Re}(\lambda_i + \delta \lambda_i) > 0\]

****
\[\delta \lambda_i = w_i^\top \delta A v_i + O(\|\delta A\|^2)\]

\textbf{ D26-1a}
\begin{itemize}
\item ****\(\min_i \text{Re}(\lambda_i) > 0\)
\item \textbf{** \(\delta \lambda_i\)   }**
\item \textbf{**\(w_i\) }**
\end{itemize}



\subsubsection{ D26-2}


****
\textbf{** \(\omega\) }**
\[H(\omega) = \|G(j\omega)\|_2 = \frac{1}{\min_i |j\omega - \lambda_i|}\]

\textbf{D26-2 }
\[\omega_{res} = \arg\max_\omega H(\omega) = \arg\max_\omega \min_i |\omega - \text{Im}(\lambda_i)|\]

\textbf{-}G-B product
\[\text{GBW} = \omega_{bw} \cdot H(\omega_{bw}) \leq \frac{\sigma_{\max}(C) \cdot \sigma_{\max}(B)}{\min_i |\text{Re}(\lambda_i)|}\]

****
\begin{itemize}
\item ****
\item \textbf{GBW } I/O
\item ****\(\min_i |\text{Re}(\lambda_i)| \to 0\)  GBW  \(\infty\)
\end{itemize}



\subsubsection{ D26-3}


****
\textbf{*}***
\[\epsilon = \inf_{i: \text{Re}(\lambda_i) > 0} \text{Re}(\lambda_i)\]

\textbf{D26-3 }
\[G_{AMP} \sim \epsilon^{-\gamma}, \quad \gamma = \arg\max_i \frac{\|\delta \lambda_i\|}{\text{Re}(\lambda_i)}\]

****
\[\tau_{relax} \sim \epsilon^{-\nu}, \quad \nu = 1 \text{}\]

\textbf{ D26-3a}
\begin{itemize}
\item \textbf{**\(\epsilon \to 0\) }**
\item \textbf{ \(\tau_{relax}\) }
\item **** +
\end{itemize}



\subsubsection{ D26-4-}


****
\textbf{-}
\[I_{causal} = \sup_{t} \mathbb{E}[y(t) | \text{do}(\delta)] - \mathbb{E}[y(t) | \delta = 0]\]

\textbf{D26-4 }
\[I_{causal} = \int_{0}^{\infty} \|h(\tau)\| \cdot \|\delta\| d\tau\]

\textbf{** \(h(\tau) = Ce^{A\tau}B\) }**

****
\[\text{Gr}_{causal} = \{(u_i \to y_j): \|CB_{ij}\| > \epsilon_{thresh}\}\]

\textbf{ D26-4a}
\begin{itemize}
\item \textbf{** } \(CB\) **
\item \textbf{** = }**
\item \textbf{*}***
\end{itemize}



\subsubsection{ D26-5}


****
****
\[\mathcal{R}(\mathcal{S}) = \sup\{\|\delta\| : \text{spec}(A+\delta A) \subseteq \text{stable region}\}\]

\textbf{D26-5 }H
\[\mathcal{R}(\mathcal{S}) = \frac{1}{\|T_{zw}(s)\|_\infty}\]

\(T_{zw}(s)\) \textbf{ \(w\)  \(z\) }

****
\[\mathcal{D}_{stable} = \{\delta : \exists \gamma > 0: \|T_{zw}(s)\|_\infty < \gamma\}\]

\textbf{ D26-5a}
\begin{itemize}
\item \textbf{ =  H }
\item \textbf{*}***
\item ****
\end{itemize}



\subsection{3.}


\subsubsection{ D26-A-}


\textbf{ D26-1  D26-3}
1. \(\ddot{x} + 2\zeta\omega_n \dot{x} + \omega_n^2 x = u\)
2.  \(\zeta\)  0.01  0.99
3.
4. ****\(\zeta \to 0\)  \(G_{AMP} \to \infty\)


\subsubsection{ D26-B}


\textbf{ D26-5}
1.  PID  \(C(s) = K_p + K_i/s + K_d s\)
2.  \(\delta K, \delta \phi\)
3. \textbf{-}
4. **** \(\mathcal{R} = 1/\|T_{zw}\|_\infty\)


\subsubsection{ D26-C}


\textbf{ D26-4}
1.
2.
3. ****
4. ****\(I_{causal}\)


\subsubsection{ D26-D}


\textbf{ D26-3}
1.
2.  \(\tau_{relax}\)
3.  \(\tau_{relax} \sim \epsilon^{-\nu}\)
4. ****\(\nu \approx 1\)


\subsection{4.}


\subsubsection{4.1 D8}


\textbf{D26-D8 }
\begin{itemize}
\item D26 \textbf{**  D8-1 }**
\item D26-3 \textbf{**  D8-2 }**
\item \textbf{** = }**
\end{itemize}



\subsubsection{4.2 D15}


\textbf{D26-D15 }
\begin{itemize}
\item D26 \textbf{** \(\mathcal{R}\)  D15-2 }**
\item D26-5 \textbf{**  D15-1 }**
\item \textbf{** = }**
\end{itemize}



\subsubsection{4.3 D1-S12}


\textbf{D26-D1-S12 }
\begin{itemize}
\item D26-3 \textbf{** = D1-S12  }\(T_c\) **
\item D26-3 \textbf{** = D1-S12  }\(d\rho/dt \to 0\)**
\item \textbf{** = }**
\end{itemize}



\subsection{5.}


\subsubsection{5.1}


1. \textbf{D26-3 }
2. \textbf{D26-5 } H
3. \textbf{D26-1 }

\subsubsection{5.2}


1. **** Lyapunov
2. ****
3. ****


\subsection{6.}


\#\#\#

**** \(\dot{x} = Ax + Bu\)  \(\delta A, \delta B, \eta\)

****
1. ****\(\delta \lambda_i = w_i^\top \delta A v_i\)
2. ****\(\omega_{res} = \arg\max_\omega \min_i |\omega - \text{Im}(\lambda_i)|\)
3. ****\(G_{AMP} \sim \epsilon^{-\gamma}\)\(\tau_{relax} \sim \epsilon^{-\nu}\)
4. ****\(I_{causal} = \int \|h(\tau)\| \|\delta\| d\tau\)
5. ****\(\mathcal{R} = 1/\|T_{zw}\|_\infty\)

\textbf{**  }\textbf{** +  + }

\textbf{*}*** \(\lambda_i\)


*D26 -\textbf{*}

\begin{thebibliography}{99}
\bibitem{ref1} Shannon, C.E. (1948). A Mathematical Theory of Communication.
\bibitem{ref2} Cover, T.M. (2006). Elements of Information Theory.
\end{thebibliography}

\end{document}